% Template:     Tesis LaTeX
% Documento:    Archivo principal
% Versión:      3.4.0 (23/08/2024)
% Codificación: UTF-8
%
% Autor: Pablo Pizarro R.
%        pablo@ppizarror.com
%
% Manual template: [https://latex.ppizarror.com/tesis]
% Licencia MIT:    [https://opensource.org/licenses/MIT]

% CREACIÓN DEL DOCUMENTO
\documentclass[
	spanish, % Idioma: spanish, english, etc.	
	letterpaper, oneside
]{book}

% INFORMACIÓN DEL DOCUMENTO
\def\documenttitle {Detección de Outliers en Curvas de Luz de Supernovas con VQ-VAE y Generación de Ejemplos Contrafactuales}
\def\documentsubtitle {}
\def\degreetitle {
	Tesis para optar al grado de magíster en ciencias de la ingeniería, mención área de la ingeniería
	\bigbreak\vspace{0.3cm}
	Memoria para optar al título de ingeniero
}

\def\universityname {Universidad de Chile}
\def\universityfaculty {Facultad de Ciencias Físicas y Matemáticas}
\def\universitydepartment {Departamento de Postgrado y Postítulo}
\def\universitydepartmentimage {departamentos/uchile2}
\def\universitydepartmentimagecfg {height=3cm}
\def\universitylocation {Santiago de Chile}

% INTEGRANTES, PROFESORES Y FECHAS
\def\documentauthor {Javiera Palominos Clarke}
\def\documentdate {\the\year}

\def\portrait {
	\begin{center}
	\vspace{1.5cm} ~ \\
	\MakeUppercase{\textbf{\documenttitle}} ~ \\
	\vspace{1.5cm}
	\MakeUppercase{\degreetitle} ~ \\
	\vfill
	\begin{tabular}{c}
		\MakeUppercase{\textbf{\documentauthor}} \\ \\
		\vspace{1.0cm} \\
		PROFESOR GUÍA: \\
		PROFESOR GUÍA \\
		\vspace{0.5cm} \\
		MIEMBROS DE LA COMISIÓN: \\
		PROFESOR 2 \\
		PROFESOR 3 \\
		\vspace{0.5cm} \\
		\MakeUppercase{\universitylocation} \\
		\MakeUppercase{\documentdate}
	\end{tabular}
	\end{center}
}

\def\abstracttable {
	\begin{tabular}{l}
		RESUMEN DE LA MEMORIA PARA OPTAR \\
		AL TÍTULO DE MAGÍSTER EN CIENCIAS \\
		DE LA INGENIERÍA \\
		POR: \MakeUppercase{\documentauthor} \\
		FECHA: \MakeUppercase{\documentdate} \\
		PROF. GUÍA: NOMBRE PROFESOR
	\end{tabular}
}

% IMPORTACIÓN DEL TEMPLATE
\input{template}

% INICIO DE LAS PÁGINAS
\begin{document}

% PORTADA
\templatePortrait

% CONFIGURACIÓN DE PÁGINA Y ENCABEZADOS
\templatePagecfg

% RESUMEN O ABSTRACT
\begin{abstractd}
	% TODO: completar resumen
	En este trabajo se presenta un pipeline de detección de anomalías y generación
	de contrafactuales para curvas de luz de supernovas del dataset ELAsTiCC-1.
	El enfoque convierte las curvas de luz multibanda en imágenes 2grid y utiliza
	un autoencoder con cuantización vectorial (VQ-VAE) con un prior tipo MaskGIT
	para detectar outliers y generar explicaciones contrafactuales interpretables.
\end{abstractd}

% DEDICATORIA
\begin{dedicatory}
	% TODO: completar dedicatoria
\end{dedicatory}

% AGRADECIMIENTOS
\begin{acknowledgments}
	% TODO: completar agradecimientos
\end{acknowledgments}

% TABLA DE CONTENIDOS - ÍNDICE
\templateIndex

% CONFIGURACIONES FINALES
\templateFinalcfg

% ======================= INICIO DEL DOCUMENTO =======================

\input{capitulos/01_introduccion}
\input{capitulos/02_marco_teorico}
\input{capitulos/03_metodologia}
\input{capitulos/04_resultados}
\input{capitulos/05_discusion_conclusiones}

% REFERENCIAS
\bibliography{library}

% ANEXOS
% \begin{appendixd}
%   \input{capitulos/anexos}
% \end{appendixd}

% FIN DEL DOCUMENTO
\end{document}
